En este capítulo se sintetizan las principales aportaciones del trabajo y se evalúa el
grado de consecución de los objetivos planteados en el Capítulo~1. La
Sección~\ref{sec:conclusiones} recoge las conclusiones del desarrollo e implementación
del sistema. La Sección~\ref{sec:lineas_futuras} identifica líneas de trabajo futuro
que podrían ampliar las capacidades de Voctomix~2.0 o adaptarla a nuevos contextos.

\section{Conclusiones}
\label{sec:conclusiones}

Este trabajo ha desarrollado y validado Voctomix~2.0, un sistema modular de producción
audiovisual en directo de código abierto que amplía la herramienta Voctomix original
con funcionalidades de nivel profesional. A continuación se evalúa el grado de
consecución de cada uno de los seis objetivos específicos planteados.

\vspace{0.3cm}
\noindent\textbf{OE-1: Funcionalidades de producción profesional.}
Se han implementado las tres funcionalidades broadcast que el sistema requería. La
rotulación dinámica permite superponer gráficos PNG y texto en tiempo real, con
fundido de entrada y salida configurable y desactivación automática en cada corte de
cámara. El gestor de estados de emisión (stream blanker) ofrece tres modos (LIVE, PAUSE y NOSTREAM) que permiten interrumpir la señal de salida sin detener el
procesamiento interno. La sincronización automática de audio con el vídeo (Audio
Follows Video) hace que el sistema ajuste los volúmenes de forma automática en cada
cambio de fuente, lo que elimina la necesidad de gestión manual del audio y reduce
la carga de trabajo del realizador durante la producción.

\vspace{0.2cm}
\noindent\textbf{OE-2: Servicio de telemetría.}
El servicio de telemetría registra en tiempo real los eventos de operación del
mezclador: el tipo CHANGE, emitido de forma inmediata al realizar cualquier acción
sobre la interfaz, y el tipo STATE, que envía el estado completo del sistema cada
5 segundos. Todos los eventos se almacenan localmente en formato JSON Lines y se
publican en un broker RabbitMQ mediante el protocolo AMQP. Esto permite integrar
el sistema con herramientas externas de monitorización o análisis sin modificar
ninguna parte del núcleo de la aplicación.

\vspace{0.2cm}
\noindent\textbf{OE-3: Contenerización con Docker.}
El sistema completo se ha empaquetado en una pila de diez contenedores Docker
orquestada con Docker Compose: voctocore, cuatro fuentes de cámara sintéticas, el
stream blanker, el gestor de audio, el broker RabbitMQ y el servicio de telemetría.
El despliegue completo se realiza con un único comando (\texttt{docker compose up}),
lo que elimina la principal barrera de adopción que suponía la instalación manual
de dependencias de GStreamer y Python.

\vspace{0.2cm}
\noindent\textbf{OE-4: Despliegue en Kubernetes.}
El sistema se ha desplegado sobre un clúster local Minikube, con manifiestos
declarativos que reproducen la misma funcionalidad que el escenario Docker Compose
pero en una infraestructura orquestada. La arquitectura del sistema escala de forma
natural desde Docker Compose hacia Kubernetes sin necesidad de modificar el código
del núcleo ni de la interfaz gráfica.

\vspace{0.2cm}
\noindent\textbf{OE-5: Validación en múltiples escenarios.}
El sistema ha sido validado funcionalmente en cuatro escenarios de despliegue: un
único PC, dos PCs en red local, Docker Compose y Kubernetes con Minikube. En los
escenarios contenerizados, la latencia de conmutación medida fue inferior a 3\,ms
en el 100\,\% de las mediciones (mediana de 1{,}5\,ms en Docker y 1{,}8\,ms en
Kubernetes), muy por debajo del umbral de un fotograma a 25\,fps (40\,ms).
El consumo de CPU con cuatro cámaras activas se mantiene por debajo del 40\,\%,
con un crecimiento sub-lineal al añadir nuevas fuentes, lo que confirma que el
sistema tiene margen de operación suficiente en equipos de gama media.

\vspace{0.2cm}
\noindent\textbf{OE-6: Despliegue en producción real.}
El sistema fue desplegado en un entorno de producción real dentro del proyecto
europeo de investigación \textbf{CyberNEMO} (UPM), donde operó como motor de
producción audiovisual integrado en la infraestructura Kubernetes de la
universidad. Este despliegue demuestra que Voctomix~2.0 va más allá del prototipo
académico y es directamente aplicable a infraestructuras de investigación y
producción audiovisual profesional.

\vspace{0.4cm}
La tabla~\ref{tab:objetivos_cumplidos} resume el grado de consecución de los
objetivos y las secciones donde se documentan.

\begin{table}[H]
  \centering
  \caption{Consecución de los objetivos específicos del TFG.}
  \label{tab:objetivos_cumplidos}
  \begin{tabular}{|c|l|c|l|}
    \hline
    \textbf{Obj.} & \textbf{Descripción} & \textbf{Estado} & \textbf{Sección} \\
    \hline
    OE-1 & Rotulación, stream blanker y AFV  & Cumplido & \S\ref{sec:overlays}, \S\ref{sec:fuentes_auxiliares} \\
    OE-2 & Telemetría JSON + RabbitMQ        & Cumplido & \S\ref{sec:telemetria} \\
    OE-3 & Contenerización Docker Compose    & Cumplido & \S\ref{sec:docker} \\
    OE-4 & Despliegue en Kubernetes          & Cumplido & \S\ref{sec:kubernetes} \\
    OE-5 & Validación multi-escenario        & Cumplido & Cap.~\ref{chap:pruebas} \\
    OE-6 & Despliegue real en CyberNEMO      & Cumplido & \S\ref{sec:escenarios_despliegue} \\
    \hline
  \end{tabular}
\end{table}

\vspace{0.3cm}

Más allá de los objetivos concretos, el trabajo ha demostrado que es posible
construir una plataforma de producción audiovisual profesional íntegramente sobre
software libre. GStreamer ha sido la elección central: su modelo de pipeline
modular ha permitido integrar cada nueva funcionalidad como un módulo independiente
sin necesidad de modificar el núcleo original, lo que facilita el mantenimiento
y la extensión del sistema por parte de terceros.

La separación entre \texttt{voctocore} como servidor multimedia y \texttt{voctogui}
como cliente ligero de control ha resultado robusta tanto en entornos contenerizados
como en red física. El protocolo de control TCP de texto plano sobre el puerto~9999
ha sido suficientemente flexible para ser consumido simultáneamente por la interfaz
gráfica, los scripts de automatización y el servicio de telemetría, sin necesidad
de mecanismos de sincronización adicionales.

En conjunto, Voctomix~2.0 es una plataforma de producción audiovisual en directo
robusta, reproducible y extensible, validada tanto en entorno de laboratorio como
en producción real, y con capacidad demostrada para operar en infraestructuras
de investigación de nivel universitario.
